\documentclass[11pt]{article}
\usepackage[margin=25mm]{geometry}
\usepackage{fontspec}
\setmainfont{Latin Modern Roman}
\usepackage{amsmath,amssymb,mathtools}
\usepackage{unicode-math}
\setmathfont{Latin Modern Math}
\usepackage{xurl,hyperref}
\hypersetup{colorlinks=true,urlcolor=blue}
\setcounter{secnumdepth}{0}
\setlength{\parindent}{0pt}
\setlength{\parskip}{5pt}
\setlength{\emergencystretch}{3em}
\newcommand{\sref}[2]{\hyperlink{source:#1:#2}{[S#2]}}
\newcommand{\cataloguescope}[1]{\paragraph{Recorded target.}#1}
\newcommand{\R}{\mathbb R}
\newcommand{\C}{\mathbb C}
\newcommand{\N}{\mathbb N}
\newcommand{\Z}{\mathbb Z}
\newcommand{\Q}{\mathbb Q}
\newcommand{\F}{\mathbb F}
\newcommand{\eps}{\varepsilon}
\begin{document}
\section{Spectral Geometry and Cyclic Homology of $SU_q(2)$}\label{statement}

\subsection{Definitions and mathematical statement}

\paragraph{Shared notation.}$\mathbb N=\{1,2,\ldots\}$; the standard number fields have their usual meanings. All additional conventions are stated locally.


\paragraph{Definitions and conventions.} For $0<q<1$, $\mathcal A(SU_q(2))$ denotes the polynomial algebra of the compact quantum-group deformation of $SU(2)$. A spectral triple $(\mathcal A,\mathcal H,D)$ represents $\mathcal A$ on a Hilbert space with a selfadjoint operator $D$, compact resolvent, and bounded commutators $[D,a]$ on the chosen algebra. The source also uses a real structure $J$ and an equivariant quantum-group action. Its ideal $K_q$ is a compact-operator ideal generated by operators whose spectra decay exponentially; the commutant and first-order axioms hold only modulo this ideal. Hochschild cycles express an orientation axiom, while cyclic cocycles encode index pairings. The metric associated to $D$ is read from the commutator seminorm $\|[D,a]\|$.

\paragraph{Source scope and qualifications.} Keep all bundled questions: Paschke duality and K\_q; orientation; a canonical local-index cocycle; differential calculus and Dirac-operator freedom; equivariant cyclic theory. The source also asks to understand a suitable smooth algebra and calculate its cyclic homology. These are distinct requirements, not one interchangeable index computation.

\paragraph{Additional source context.} The report first asks for a satisfactory $C^\infty$ completion of the polynomial algebra and its cyclic homology. In the local-index question the desired geometric three-cocycle is distinguished from the already computed one-cocycle image.

\paragraph{Catalogue mathematical target.}
\paragraph{Catalogue wording retained for traceability.}
For the spectral geometry of $SU_q(2)$, what are the geometric roles of Paschke duality and the ideal $K_q$? Can the Hochschild-cycle axiom be satisfied or suitably modified? Is there a canonical local-index cocycle from the spectral data? What differential calculus and metric arise, and what freedom in the Dirac operator changes that metric? What is the correct equivariant cyclic-homology framework, and how can it be computed?

\paragraph{Saved literature context.} The supplied catalogue's assessment dated 2026-08-21 records: This is a bundled list of open questions in q-geometry concerning Hochschild cycles, local index data, calculi, metrics, and equivariant cyclic homology; no single general resolution was located. The assessment is retained as dated provenance; it is not a new comprehensive literature or open-status certification.
\paragraph{Scope and source qualification.}Keep all bundled questions: Paschke duality and K_q; orientation; a canonical local-index cocycle; differential calculus and Dirac-operator freedom; equivariant cyclic theory. The source also asks to understand a suitable smooth algebra and calculate its cyclic homology. These are distinct requirements, not one interchangeable index computation.
\subsection{Short English statement}
Complete the geometric, index, differential, metric, and equivariant homology interpretation of the spectral data of the deformed special unitary group.
\subsection{Research attempt}
\paragraph{Model, date and investigated scope.}
OpenAI GPT-6 Astra (configured model \texttt{gpt-6-astra}), ultra reasoning; generated 11 October 2026. Proves exact metric scaling and a sufficient class of metric-invisible bounded perturbations for the supplied spectral-triple setting.

\paragraph{Sources and literature check.}
The complete supplied definition, including source qualifications and the dated literature assessment, was reviewed. The calculation below is an elementary restricted attempt motivated by [S2]; it does not claim novelty or a fresh literature-status determination.

\paragraph{Context and investigated scope.} The complete supplied definition, including source qualifications and the dated literature assessment, was reviewed. The calculation below is an elementary restricted attempt motivated by [S2]; it does not claim novelty or a fresh literature-status determination.

\paragraph{Metric freedom that survives the axioms.} Let $(\mathcal A,\mathcal H,D)$ be any one of the supplied spectral triples, and define its extended state distance by
\[
 d_D(\varphi,\psi)=\sup\{|\varphi(a)-\psi(a)|:a=a^*,\ \|[D,a]\|\leq1\}.
\]
For $c>0$, $D_c=cD$ has compact resolvent whenever $D$ does, and $[D_c,a]=c[D,a]$ is bounded. Replacing $a$ by $a/c$ in the supremum proves $d_{D_c}=c^{-1}d_D$, including infinite distances. Equivariance is unchanged. Any first-order relation holding modulo the fixed linear ideal $K_q$ still holds, because its error is multiplied by $c$.

\paragraph{An invisible perturbation.} If $B=B^*$ is bounded and commutes with every $a\in\mathcal A$, then $[D+B,a]=[D,a]$, so the metric is identical. Compact resolvent is retained: for a nonreal $z$ with $\|B(D-z)^{-1}\|<1$, the resolvent identity expresses $(D+B-z)^{-1}$ as a compact resolvent times a bounded inverse. The inverse exists by the convergent geometric series. This does not guarantee that an arbitrary $B$ preserves reality or equivariance; those must be separately imposed.

\paragraph{Obstruction to a canonical answer.} Even before analyzing $SU_q(2)$ specifically, the listed spectral-triple axioms admit an overall metric scale, while some bounded changes are metrically invisible. A ``canonical metric'' therefore requires additional normalization and a classification of permitted noncommuting perturbations. No smooth completion, Hochschild orientation cycle, geometric three-cocycle, Paschke-duality interpretation of $K_q$, differential calculus, or equivariant cyclic homology computation is produced. The metric calculation cannot substitute for those separately bundled requirements.

\paragraph{Verification and adversarial audit.}
Checked infinite distances, compact resolvent under scaling and bounded perturbations, and which reality/equivariance properties require extra hypotheses. Model self-audit only; no independent human proof checking. Prompts A and B were applied to this canonical record and its qualified source scope. Checked infinite distances, compact resolvent under scaling and bounded perturbations, and which reality/equivariance properties require extra hypotheses. Retained unresolved status and zero original-target progress; Additional normalization and control of permitted Dirac perturbations are absent, as are the other bundled homological constructions.

Coordinator checked metric scaling and commuting bounded perturbations, including the compact-resolvent factorization and separately required symmetry conditions. This additional check was model-assisted, not independent human verification.

\paragraph{Tools, process and contributions.}
Python 3 was used for file preparation and structural checks; no numerical experiment or theorem prover was used. Contribution prompt A was followed by prompt B. GPT-6 Astra with ultra reasoning drafted the research and performed its stated self-audit; any additional coordinator contribution is identified in the verification and audit. A GPT-6-based Codex coordinator assembled the manuscript and checked repository integration; its exact variant and reasoning configuration were not exposed. No human mathematical contribution, independent human proof check or formal proof certification is supplied. Complete derivations and the consulted source references are included here; private reasoning transcripts are not published.

\paragraph{Final status and first remaining gap.}
\textbf{UNRESOLVED}. Additional normalization and control of permitted Dirac perturbations are absent, as are the other bundled homological constructions.

\section{Completion Estimate}
\noindent\textbf{COMPLETION: 0\%}
The calculation is an auxiliary or restricted attempt and establishes no newly verified portion of the full source target. Additional normalization and control of permitted Dirac perturbations are absent, as are the other bundled homological constructions.

\subsection{Sources}
\begin{itemize}
\item[\hypertarget{source:statement:1}{S1}] ULAM.AI and the UnsolvedMath Contributors, UnsolvedMath, catalogue entry. Catalogue curation under CC BY 4.0; definitions and English exposition added locally.
\url{https://huggingface.co/datasets/ulamai/UnsolvedMath}.
\item[\hypertarget{source:statement:2}{S2}] Noncommutative Geometry and Quantum Field Theory (2005), Oberwolfach Reports, DOI 10.4171/owr/2005/48, pp. 2705--2760; official report source and mathematical context.
\url{https://doi.org/10.4171/owr/2005/48}.
\item[\hypertarget{source:statement:3}{S3}] Official EMS Press PDF of the same Oberwolfach report; PDF pages 44--45, printed pages 2748--2749.
\url{https://ems.press/content/serial-article-files/46022}.
\end{itemize}
\label{end:statement}
\end{document}
